\section{Antecedentes}
\label{Antecedentes}

Los drones comerciales comenzaron a utilizarse a gran escala en el conflicto Ruso-Ucraniano en el año 2022, marcando un punto de inflexión en el uso de tecnología accesible y versátil en el campo de batalla. Este fenómeno no solo subraya la capacidad de los drones para adaptarse rápidamente a las necesidades militares, sino que también refleja una tendencia histórica en la evolución de los sistemas de armas y defensa.

La historia de la guerra es, en esencia, una constante evolución entre las herramientas ofensivas y los mecanismos defensivos. Durante la Edad Media, los castillos representaban una defensa formidable contra los ejércitos a pie, proporcionando refugio y ventajas estratégicas. Sin embargo, la invención de la pólvora y el desarrollo de cañones y armas de fuego hicieron obsoletos tanto a los castillos como a las armaduras tradicionales, redefiniendo las estrategias de combate.

De manera similar, los avances en los misiles antiaéreos impulsaron la evolución de los aviones de combate, dando lugar a modelos como el SR-71 Blackbird, diseñado para superar los sistemas de defensa convencionales mediante su velocidad y altitud extremas. Ahora, los drones han emergido como una alternativa altamente económica y fácil de desplegar, con la capacidad de infligir daño localizado tanto a personal como a vehículos. Esta nueva generación de tecnología bélica permite a los combatientes alcanzar objetivos estratégicos con precisión y a un costo significativamente menor en comparación con los sistemas tradicionales.

El conflicto en Ucrania ha puesto de manifiesto estas tendencias, demostrando cómo los drones pueden ser utilizados tanto para vigilancia como para ataques precisos, consolidando su posición como una herramienta clave en los conflictos modernos.